\section{Динамика без трения}
\subsection{Законы Ньютона}
%%%%%%%%%%%%%%%%%%%%%%%%%%%%%%%%%%%%%%%%%%%%%%%%%%%%
%%%%%%%%%%%%%%%%%%%%%%%%%%%%%%%%%%%%%%%%%%%%%%%%%%%%


 \z{Вагон движется в горизонтальной плоскости со скоростью $v$ по закруглению радиусом $R$. Во сколько раз возрастёт вес груза в движущемся вагоне по сравнению с весом груза в неподвижном вагоне?}
%
{$\alpha = \sqrt{1+\frac{v^4}{R^2g^2}}$}
%%%%%%%%%%%%%%%%%%%%%%%%%%%%%%%%%%%%%%%%%%%%%%%%%%%%
\z[{\begin{wrapfigure}[4]{r}{0.3\textwidth}%
		\vspace{-12pt}%
		\begin{flushright}%
 			\includegraphics[width=0.3\textwidth]{4.1.2}%
		\end{flushright}%
	\end{wrapfigure}%
	}]%
	{Четырьмя натянутыми нитями груз закреплен на тележке. Сила натяжения горизонтальных нитей соответственно $T_1$ и $T_2$, а вертикальных --- $T_3$ и $T_4$. С каким ускорением тележка движется по горизонтальной плоскости?}
%
{$a=\frac{T_2-T_1}{T_4-T_3}g$}
%%%%%%%%%%%%%%%%%%%%%%%%%%%%%%%%%%%%%%%%%%%%%%%%%%%%	
 \z[{\begin{wrapfigure}[5]{r}{0.3\textwidth}%
		\vspace{-12pt}%
		\begin{flushright}%
 			\includegraphics[width=0.3\textwidth]{4.1.3}%
		\end{flushright}%
	\end{wrapfigure}%
	}]%
	{Верёвка выдерживает груз максимальной массы $m_1$ при его движении с некоторым ускорением, направленным вверх, и груз максимальной массы $m_2$ при его движении с ускорением, в $n$ раз меньшим первого по величине и направленным вниз (см. рисунок). Груз какой максимальной массы можно подвесить к верёвке в покое?}
 %
{$m=\frac{m_1m_2(n+1)}{nm_1+m_2}$}
%%%%%%%%%%%%%%%%%%%%%%%%%%%%%%%%%%%%%%%%%%%%%%%%%%%% 
 \z{В вертикальном цилиндре неподвижно удерживают два тонких массивных поршня: верхний -- массой $m$, нижний -- массой $2m$. Нижний поршень находится на небольшой высоте $H$ от дна цилиндра, расстояние между поршнями равно $L>>H$. Выше и ниже поршней вакуум, а между ними содержится газ. Поршни одновременно отпускают. Нижний поршень упал на дно цилиндра через время $t$ после отпускания. На какой высоте над дном находился верхний поршень в момент удара нижнего о пол? Ускорение свободного падения $g$, трения нет, массой газа по сравнению с массой поршней и изменением давления газа при движении поршней можно пренебречь.}
%
{$h=3H+L-\frac{3}{2}gt^2$}
%%%%%%%%%%%%%%%%%%%%%%%%%%%%%%%%%%%%%%%%%%%%%%%%%%%% 
 \z{Сидевший на корточках человек резко выпрямляется и, оттолкнувшись от пола, подпрыгивает так, что его центр масс поднимается на высоту $h$, равную $3/4$ его роста $l$ (высота отсчитывается от пола). Найдите среднюю силу, с которой человек действует на пол во время отталкивания. Центр масс человека, когда он стоит выпрямившись, находится на высоте $l/2$ от пола. Перед прыжком центр масс человека находился на высоте $l/4$ от пола. Масса человека $m = 75\,\text{кг}$.}
%
{$N=2mg=1500\,\text{Н}$}
%%%%%%%%%%%%%%%%%%%%%%%%%%%%%%%%%%%%%%%%%%%%%%%%%%%%